% Accepted and integrated as sections/A_F4_all_genus_two.tex after checkpoint10.
\subsection{Every genus-two curve over \texorpdfstring{$\F_4$}{F4}}

\begin{claim}{A.C2.26}{Proved}
Every smooth proper geometrically connected
genus-two curve over $\F_4$ satisfies A1.
The conclusion holds over every finite
extension, for every finitely generated
form of any of these geometric curves,
and for every hyperbolic open subcurve.

The twelve geometric two-rank-one classes
fixed by fourth-power Frobenius are
\[
 C(A,B):\quad y^2+y=Ax^3+Bx+\frac1x,
 \qquad A\in\F_4^*,\quad B\in\F_4.                         \tag{C2.92}
\]
Explicit odd theta-torsion annihilators
$5,7,11,13,19,39$ settle all of them.
Together with A.C2.25 and the two-rank-zero
argument, this covers all $64$ geometric
genus-two moduli classes fixed by
fourth-power Frobenius. Universal A1
over arbitrary finite fields remains A.30.
\end{claim}
\begin{proof}
\proofparagraph{The exact twelve-class normalization.}
Over an algebraically closed field of
characteristic two, the genus-two
hyperelliptic quotient is separable and
admits a reduced Artin--Schreier equation.
For odd pole orders $m_j$, the genus and
two-rank formulas are
\[
 2g-2=-4+\sum_j(m_j+1),\qquad
 f(C)=\#\{\text{poles}\}-1
\]
\citep[Notation 1.1, author p.~2]{A_ElkinPries2013}.
When $g=2$ and $f(C)=1$, there are two
poles, of orders three and one. Send
them to infinity and zero. Eliminating
even terms and an additive constant gives
$ax^3+bx+c/x$ with $ac\ne0$; a unique
scaling normalizes $c=1$.
These are precisely the case $(1,3)$
normal forms of
\citet[Proposition 14, author p.~14]
{A_CardonaNartPujolas2005}.

The normalized pair $(A,B)$ is a complete
geometric invariant. An isomorphism
preserves the unique hyperelliptic pencil.
The branch points have different different
exponents, four and two, so the induced
base map fixes both and is $x\mapsto ux$.
A generator change is $y\mapsto y+v(x)$.
A nonconstant $v$ gives an even leading
pole order in $v^2+v$, whereas the difference
of the reduced equations has only odd
pole orders. Thus $v$ is constant.
The simple-pole coefficient forces $u=1$,
and then $A,B$ agree. This is also the
scaling criterion in the cited proposition.

Fourth-power Frobenius fixes the class
exactly when $A^4=A$ and $B^4=B$.
There are $3\cdot4=12$ possibilities,
all represented by actual equations
over $\F_4$. Under coefficientwise
squaring the seven representatives are
\[
 (1,0),(1,1),(1,\alpha),(\alpha,0),
 (\alpha,1),(\alpha,\alpha),(\alpha,\alpha^2),
 \qquad \alpha^2+\alpha+1=0.                               \tag{C2.93}
\]
The first two are fixed, and the others
have distinct conjugates. No descent
assertion from a field of moduli is
required for this count.

\proofparagraph{The theta origin and a norm criterion.}
Let $W_0,W_\infty$ lie above the branch
values. The valuations are
\[
 \begin{array}{c|rr}
       &x&y\\\hline
 W_0&2&-1\\
 W_\infty&-2&-3
 \end{array}
\]
and the different is $2W_0+4W_\infty$.
Consequently
\[
 \operatorname{div}(dx)=2W_0,
 \qquad B_\Theta=\mathcal O(W_0),
 \qquad \Theta=\{[R-W_0]:R\in C\}.                         \tag{C2.94}
\]
Here $B_\Theta$ is the line bundle of
locally exact differentials. In particular
zero is on theta, so a nonzero witness
must be checked. In the polynomial model
$Y=xy$, one has
\[
                 Y^2+xY=Ax^5+Bx^3+x.
\]
The pair $(x,0)$ represents the two-torsion
class $[W_0-W_\infty]$, so adding it to
$[R-W_\infty]$ gives the correct theta
class $[R-W_0]$. The fixed norm proof
below is independent of the discovery
calculation with this Jacobian model.

Let $n>1$ be odd and $d=(n-1)/2$.
Suppose a finite extension contains
$t\ne0,y_0$, polynomials $a,b$ and a
nonzero scalar $c$ such that
\[
 \begin{aligned}
 &\deg a\le d,\qquad \deg b\le d-2,\\
 &x(a^2+ab)+b^2(Ax^4+Bx^2+1)=c(x+t)^n,\\
 &a(t)+b(t)y_0=0,\qquad b(0)b(t)\ne0,\\
 &y_0^2+y_0=At^3+Bt+1/t.
 \end{aligned}                                             \tag{C2.95}
\]
Put $R=(t,y_0)$ and
\[
                            q=\frac{a+by}{x^d}.
\]
At $W_0$, the term $b(0)y$ has valuation
$-1$, strictly smaller than that of $a$;
the denominator has valuation $2d$.
Thus the pole order is exactly $n$.
At infinity, $a/x^d$ has nonnegative
valuation and $by/x^d$ has valuation at
least one. There is no other pole.
Equation (C2.95) says
\[
                  q\,\iota(q)=c\frac{(x+t)^n}{x^n}.
\]
Every zero therefore lies above $x=t$.
At $R$ its numerator is zero, while at
the conjugate point it equals $b(t)\ne0$.
This fibre is unramified, so all $n$
zeros lie at $R$. We have proved the
exact divisor identity
\[
                      \operatorname{div}(q)=n(R-W_0).      \tag{C2.96}
\]
The class $[R-W_0]$ is nonzero: otherwise
a degree-one function would make the
genus-two curve rational. Its actual odd
order $e>1$ divides $n$; no proof that
$e=n$ is needed.

\proofparagraph{Seven fixed polynomial identities.}
The following table specifies every
representative in (C2.93). Integer field
elements are binary polynomial coefficients
in ascending powers of a generator $w$.
The column $m$ is the binary modulus;
$g$ is a checked primitive multiplicative
generator. The columns $\alpha,t,y_0$
use that same encoding.
\begin{center}
\begin{tabular}{crrrrrr}
\toprule
$(A,B)$&$|F|$&$m$&$g$&$\alpha$&$(t,y_0)$&$n$\\
\midrule
$(1,0)$&16&19&2&6&$(6,2)$&13\\
$(1,1)$&4&7&2&2&$(1,2)$&5\\
$(1,\alpha)$&16&19&2&6&$(7,2)$&11\\
$(\alpha,0)$&256&283&3&189&$(80,242)$&39\\
$(\alpha,1)$&16&19&2&6&$(6,4)$&19\\
$(\alpha,\alpha)$&4&7&2&2&$(1,2)$&7\\
$(\alpha,\alpha^2)$&16&19&2&6&$(7,2)$&5\\
\bottomrule
\end{tabular}
\end{center}
The three moduli are
$w^2+w+1$, $w^4+w+1$, and
$w^8+w^4+w^3+w+1$.
The complete coefficients of $a,b,c$
for all seven identities are in
\path{code/A_F4_rank_one_witnesses.json}.
For example the degree-seven row has
\[
 a=\alpha^2+x+\alpha x^2+x^3,
 \qquad b=1+\alpha x,
 \qquad c=1.
\]
The largest row has degree $39$ and
uses only $\F_{256}$.

The fixed data have SHA256
\begin{center}
\path{82dd18cfcfac963e1a4217e03a55b39a47766ee521c6601f3704e9828b417ec3}.
\end{center}
The standalone direct verifier
\path{code/A_F4_rank_one_norm_checks.py}
and its JSON output check every coefficient
of (C2.95), the point and local conditions,
complete generator cycles, and the
twelve-class conjugation coverage.
It imports no C arithmetic, search,
jet solver or Jacobian implementation.
The independent verifier
\path{code/A_F4_rank_one_degeneration.py}
and its output use freshly implemented
binary-field and polynomial arithmetic;
all seven identities pass there too.
These are positive fixed norm certificates,
independent of any exhaustive negative scan.

\proofparagraph{Primitive vanishing and all three strata.}
Direct multiplication gives
\[
                 \langle2,-1\rangle=(\Z/n)^*
       \qquad(n=5,7,11,13,19,39).                           \tag{C2.97}
\]
At $39$ the doubling order is $12$ and
negation supplies the second coset, giving
all $24$ units. Reduction of units from
$n$ to any divisor $e$ is surjective, so
(C2.97) holds at the actual order of the
nonzero class $\eta=[R-W_0]$ from (C2.96).
Theta is symmetric because
$R+\iota(R)\sim2W_0$. Thus $\eta$ and
$-\eta$ supply a theta point on every
primitive doubling orbit at $e$.
By A.C2.5 the connected cyclic \'etale
cover of actual degree $e$ has primitive
factor of two-rank zero. A.C2.19 or
A.C2.20 proves A1. There is no extraction
of a primary component that might leave
theta, and no claim about the survival
of lower-conductor characters is needed.

Squaring coefficients gives the five
conjugate rows and all twelve classes.
Finite-extension descent, finitely
generated forms and open subcurves follow
from the geometric criterion, A.3 and
A.8, as in A.C2.25. In particular all
isomorphism and cover data are first
defined after a finite extension of
the whole base field.

Every genus-two curve has two-rank
zero, one or two. Rank one is settled
above, and rank two is A.C2.25. At rank
zero, $J[2](\overline k)=0$, and repeated
doubling excludes every nonzero geometric
two-primary torsion point. A.1 turns a
section collision into a two-primary
point difference, which is therefore
zero. The Abel map is injective, giving
the conclusion in that stratum over
every finitely generated characteristic-two
field.

For the count of geometric classes,
the invariant description of
\citet[Theorem 20, author p.~19;
Theorem 22 and Corollary 23, author
p.~21]{A_CardonaNartPujolas2005}
identifies the moduli with three invariant
coordinates. Over $\F_4$ the three
strata have respectively $4$, $12$, and
$48$ geometric classes, for a total of
$64$. This distinguishes geometric
classes from their arithmetic twists;
all twists are included in the forms
assertion. The cited classification
does not assert the prime-to-two section
theorem proved here by the certificates.

The full independent review is
\path{audit/A_F4_rank_one_degeneration_review.md}.
The earlier theta-five and theta-thirteen
results remain valid and have direct
witnesses in the table. The conclusion
is the stated finite-field special case;
no universal existence of a suitable
primitive-zero cover is inferred.
\end{proof}
